% !Mode:: "TeX:UTF-8"

\chapter{绪论}

\section{研究背景}

大语言模型在常识问答、代码生成和文本分析等任务中表现突出，但其是否真正具备稳定的空间建模能力仍需要通过可控实验验证。空间推理问题通常同时包含自然语言理解、几何约束识别、数量关系比较和逻辑一致性保持等要求，因此适合作为分析模型认知边界的实验载体。

本项目选择“电动轮椅在狭窄楼梯转角处能否顺利完成 90 度转弯”作为固定任务。该问题具有现实空间约束，同时又可以被改写为自然语言描述、参数化描述和形式化建模描述三种层级。与直接求解几何可通行性不同，本报告关注的是模型如何回答、为何出错，以及哪些描述变化会诱发回答翻转。

\section{研究目标}

本报告拟完成以下目标：

\begin{enumerate}
    \item 构建包含描述层级、干扰条件和场景临界性的 Prompt 实验矩阵；
    \item 自动化采集不同模型在同一空间任务上的回复；
    \item 从回复文本中抽取可量化行为特征；
    \item 使用数据挖掘方法识别模型的准确率差异、逻辑一致性、噪声敏感性和典型失效模式；
    \item 为课程报告和答辩提供可复现的数据、图表和结论框架。
\end{enumerate}

\section{项目技术路线}

项目整体流程如图 \ref{fig:pipeline} 所示：首先生成多层级 Prompt，然后调用 mock 或真实模型 API 采集回复；随后从回复中抽取最终判断、公式使用、坐标建模和错误类别等特征；最后进行统计分析、聚类分析和关联规则挖掘。

\begin{figure}[htbp]
    \centering
    \fbox{\parbox{0.86\textwidth}{
    Prompt 生成 $\rightarrow$ 模型回复采集 $\rightarrow$ 行为特征抽取 $\rightarrow$ 统计分析 / 聚类 / 关联规则 $\rightarrow$ 报告与可视化
    }}
    \bicaption{项目数据挖掘流程}{Data mining workflow of the project}
    \label{fig:pipeline}
\end{figure}

\section{本文组织结构}

第 2 章介绍数据集构造与实验变量设计；第 3 章说明模型回复采集和行为特征工程；第 4 章介绍准确率统计、聚类、决策树和 FP-Growth 等挖掘方法；第 5 章给出实验结果分析框架；最后总结当前工作并讨论后续完善方向。
